\section{A rational family and finite central-neighborhood sequences}
\label{sec:discrete}
We now separate central and arbitrary movement on the same explicit
sparse linear program. In the unit-scale results preceding the final
subsection, $r\to\infty$ and $F(x)=\sum_{i=1}^r b(x_i)$ on $(-1,1)^r$.
The objective is to minimize $-w^\top x$, equivalently to maximize the
support direction $c=w$, so $g_c(x)=\sum_iw_i(1-x_i)$.
Sign changes are barrier isometries, so arbitrary prescribed objective
signs give the same results.

Define the following exact rational and integer data:
\begin{equation}
 h_j=\frac1{j\lceil\sqrt j\rceil},\quad
 A_i=\left\lfloor4r\sum_{j<i}h_j\right\rfloor,\quad
 a_i=A_i\log2,\quad w_i=2^{-A_i},\quad
 T=a_r,\quad\eps_r=r2^{-A_r}.
 \label{eq:dyadic-family}
\end{equation}
The box has two inequalities per coordinate, with one nonzero per row
and two per column. Since $A_i=\Theta(r)$ for every $i\geq2$, the
ordinary rational numerator--denominator encoding of the objective has
$\Theta(r^2)$ bits, and $\eps_r$ has $\Theta(r)$ bits. Integer square
roots, rational summation and flooring construct the data in polynomial
time. An exponent encoding is a different, more succinct representation.

\begin{theorem}[Three primal comparisons on the dyadic family]
\label{thm:dyadic}
Let $x(s)$ denote the exact center at $\eta=e^s$. For all sufficiently
large $r$, its first $\eps_r$-accurate logarithmic parameter satisfies
$T-1<s_\eps\leq T$, and
\begin{equation}
 L_\opt(\eps_r)=\Theta(r\sqrt{\log r}),\quad
 d_F(0,x(s_\eps))=\Theta(r\sqrt{\log r}),\quad
 L_\CP(\eps_r)=\Theta(r\log r).
 \label{eq:dyadic-lengths}
\end{equation}
Consequently, for fixed $0<R<1$, the smallest number of arbitrary
forward-$R$ moves to accuracy is $\Theta_R(r\sqrt{\log r})$.
For the equal-scale standard barriers in rank $r$, the supremum of
$L_\CP(\eps)/L_\opt(\eps)$ is $\Theta(\sqrt{\log r})$.
\end{theorem}
\begin{proof}
Write $\bar a_i=4r(\log2)\sum_{j<i}h_j$ and
$e_i=a_i-\bar a_i\in(-\log2,0]$. The inequalities
$\frac12j^{-3/2}\leq h_j\leq j^{-3/2}$ show $T=\Theta(r)$ and
\begin{equation}
 T-a_i=4r(\log2)\sum_{j=i}^{r-1}h_j+O(1),\quad
 a_{j+1}-a_j\geq2r(\log2)j^{-3/2}-\log2.
 \label{eq:dyadic-tail}
\end{equation}
The upper error estimate in \cref{prop:distribution-law} gives
$E(e^T)\leq re^{-T}$. For a sufficiently large absolute constant $K$,
all $i\leq r-K\sqrt r$ have $T-a_i\geq1$: sum the lower bound on
$h_j$ over the last $K\sqrt r$ terms and use the bounded rounding error.
At $s=T-1$, these channels have $e^{s-a_i}\geq1$, so
$E(e^{T-1})\geq e^{-(T-1)}(2-\sqrt2)(r-O(\sqrt r))>re^{-T}$.
Strictly decreasing error proves the asserted window.

For $u\geq-1$, \eqref{eq:profile-tails} gives
$H(u)=\Theta(1+u_+)$ with universal constants. Tail summation in
\eqref{eq:dyadic-tail} gives
$\sum_i(T-a_i)^2=\Theta(r^2\log r)$: the upper bound follows from
$T-a_i\leq Cr/\sqrt i$, and the lower from
$T-a_i\geq cr/\sqrt i$ for $i\leq r/16$, after adjusting constants
and discarding finitely many small $r$. The accurate-center window
therefore proves the middle assertion of \eqref{eq:dyadic-lengths} and
an upper bound on $L_\opt$.
For the lower bound on $L_\opt$, every feasible allocation satisfies
$w_i(1-x_i)\leq\eps_r$, hence
$\rho(x_i)\geq[T-a_i-\log r]_+$. On $i\leq\sqrt r$ these lower
bounds are at least $c'r/\sqrt i$ for large $r$, and their squared sum
is $\Omega(r^2\log r)$. Apply \cref{thm:water-filling}.

Let $J=\lfloor r^{2/3}\rfloor$. On $[a_j,a_{j+1}]$ the first $j$
velocities are at least $v_0$. These intervals for $j\leq J$ end before
$s_\eps$, since $T-a_{J+1}=\Theta(r^{2/3})$. Their total length is at
least $v_0\sum_{j\leq J}\sqrt j(a_{j+1}-a_j)$. Summation by parts
bounds the rounding contribution
$\sum_{j\leq J}\sqrt j(e_{j+1}-e_j)$ in absolute value by
$2(\log2)\sqrt J$. The unrounded sum is
$4r(\log2)\sum_{j\leq J}\sqrt jh_j=\Theta(r\log r)$.
This proves the central lower bound; \cref{thm:sharp-prefix} gives
its matching upper bound. \Cref{cor:water-movement} gives the optimal
count, and \cref{thm:scalar-dilation} gives the universal same-accuracy
upper ratio. Scalar intervals embed as diagonal spectral or Jordan
channels, so the lower ratio applies throughout the claimed class.
\end{proof}

\subsection{Arbitrary labels, actual accuracy and a growing metric tube}
For a finite sequence $z_0,\ldots,z_N$, the following result requires the
actual initial point and objective accuracy, rather than prescribed
starting and ending central parameters. Reference labels may backtrack.

\begin{theorem}[Discrete neighborhood lower bound]
\label{thm:tube}
Fix $0<R<1$. For the family \eqref{eq:dyadic-family}, suppose
\begin{equation}
 z_0=0,\quad g_c(z_N)\leq\eps_r,\quad
 \norm{z_{k+1}-z_k}_{F,z_k}\leq R,\quad
 d_F(z_k,x(s_k))\leq\delta_r
 \label{eq:tube-contract}
\end{equation}
for arbitrary labels $s_k\in\R\cup\{-\infty\}$, where
$x(-\infty)=0$. If $\delta_r=o(r^{2/3})$ and
$C_r=2\delta_r-\log(1-R)$, then
\begin{equation}
 N=\Omega_R\left(\frac{r\log r}{C_r\sqrt{1+C_r}}\right).
 \label{eq:growing-tube}
\end{equation}
In particular, a fixed tube has lower count $\Omega_{R,\delta}(r\log r)$,
matching exact central-arc sampling. Relative to optimal unrestricted
movement, the lower ratio diverges whenever
$\delta_r=o((\log r)^{1/3})$.
\end{theorem}
\begin{proof}
The triangle inequality and \cref{lem:chord-distance} bound consecutive
center distances by $C_r$. Define
\[
 J=\lfloor r^{2/3}\rfloor,\quad b=a_{J+1},\quad
 n(u)=\#\{i:a_i\leq u\},\quad
 P(s)=\int_0^{\min(\max(s,0),b)}\sqrt{n(u)}\dd u.
\]
For large $r$, the first $J$ threshold gaps are at least $g=\log2$ by
\eqref{eq:dyadic-tail}. Also $P(b)=\Theta(r\log r)$ by the harmonic
sum in the proof of \cref{thm:dyadic}.

Consider any pair of labels and clip both into $[0,b]$. Because every
$H(s-a_i)$ increases, clipping cannot increase the Euclidean distance
between their coordinate vectors. Order the clipped labels as $u\leq t$.
If $u=t$ there is no progress change. Otherwise put $j=n(u)\geq1$,
$\Delta=t-u$, and $D=C_r/v_0$. The first $j$ coordinates have velocity
at least $v_0$ throughout $[u,t]$, so $\Delta\leq D/\sqrt j$.
Uniform threshold spacing gives $n(t)\leq j+1+\Delta/g$ (endpoint
values do not affect the integral). Therefore
\begin{equation}
 |P(s_{k+1})-P(s_k)|
 \leq\Delta\sqrt{j+1+\Delta/g}
 \leq D\sqrt{2+D/g}=O(C_r\sqrt{1+C_r}).
 \label{eq:one-progress-jump}
\end{equation}
This argument includes jumps across either clipping boundary and
negative label moves.

We next force the required endpoint progress. Actual accuracy and
$w_1=1$ imply $z_{N,1}\geq1-\eps_r$, since each coordinate's objective
gap is nonnegative. Scalar metric contraction and the terminal tube give
$H(s_N)\geq\rho(1-\eps_r)-\delta_r\geq T-\log r-\delta_r$.
Since $H(s)\leq1/\sqrt2+s_+$, this implies
$s_N\geq T-\log r-\delta_r-1/\sqrt2>b$ for large $r$;
here $T-b=\Theta(r^{2/3})$ is essential.
For the initial progress we can bound the entire clipped step-profile
arc by the smooth central arc: on $[0,b]$ its speed $\sqrt{n(s)}$
is at most $\norm{(v(s-a_i))}_2/v_0$. Hence
\[
 P(s_0)\leq v_0^{-1}L_\CP(e^{s_0})
 \leq v_0^{-1}\Gamma_r d_F(0,x(s_0))
 \leq v_0^{-1}\Gamma_r\delta_r=o(r\log r).
\]
For $s_0\leq0$ the first inequality is immediate because $P(s_0)=0$;
for $s_0>0$ it follows by integration. Thus the net progress is
$\Theta(r\log r)$. Sum absolute increments in
\eqref{eq:one-progress-jump}; no monotonicity hypothesis is used.
The final ratio follows from \cref{thm:dyadic} because
$C_r\sqrt{1+C_r}=o(\sqrt{\log r})$ in the stated regime.
\end{proof}
The label $-\infty$ simply includes the analytic center in the exact
central sequence, even for tube radius zero. If finite labels are required,
the same proof holds; an exact radius-zero tube with analytic-center
initialization would then be an empty class.

\begin{lemma}[Newton decrement implies a metric tube]
\label{lem:decrement-tube}
Let $f=F-\ip{c}{\cdot}$ have an interior minimizer $x^\star$, where $F$
is standard self-concordant. Write
$\lambda_f(z)=\norm{\nabla f(z)}_{F,z,*}$. If
$\lambda_f(z)\leq\beta<1/2$, with $\beta\geq0$, then
\begin{equation}
 d_F(z,x^\star)\leq
 \delta_\beta:=\log\frac{1-\beta}{1-2\beta}.
 \label{eq:decrement-tube}
\end{equation}
\end{lemma}
\begin{proof}
Put $h=x^\star-z$ and $t=\norm{h}_{F,z}$. Along the segment, the scalar
local norm $a(u)=\norm{h}_{F,z+uh}$ satisfies $|a'|\leq a^2$,
hence $a(u)\geq t/(1+ut)$. Since $\nabla f(x^\star)=0$,
\[
 -Df(z)[h]=\int_0^1D^2f(z+uh)[h,h]\dd u
 \geq\int_0^1\frac{t^2}{(1+ut)^2}\dd u=\frac{t^2}{1+t}.
\]
Dual Cauchy--Schwarz gives $t\leq\lambda_f(z)/(1-\lambda_f(z))<1$.
The straight-segment upper bound is $-\log(1-t)$; substituting $\beta$
proves the claim. If $t=0$, the conclusion holds directly.
\end{proof}
Therefore \cref{thm:tube} applies to every forward-$R$ sequence that
maintains the usual fixed decrement condition
$\norm{\nabla F(z_k)-e^{s_k}w}_{F,z_k,*}\leq\beta<1/2$ and returns
an actually accurate point. It also allows $\beta_r\uparrow1/2$ whenever
$\delta_{\beta_r}=o((\log r)^{1/3})$ for a diverging overhead.

There is a useful precise transfer to feasible conic methods. Suppose a
cone barrier $\widehat F$ restricts on an affine slice $z=z_0+Bx$ to
$F(x)$ up to an additive constant, and
suppose $q=c-A^*y$ is a feasible slack for a conic minimization problem,
with $AB=0$ and $B^*c=-w$. The dual norm of the restricted
covector $B^*\ell$ in the restricted Hessian is at most the full dual
norm of $\ell$: both are suprema of $\ell[h]$ over local unit vectors,
and restriction reduces that set. Consequently the conventional residual
condition
\[
 \norm{q+\mu\nabla\widehat F(z)}_{\widehat F,z,*}\leq\beta\mu
\]
implies the primal decrement bound at parameter $1/\mu$, since
$B^*(q+\mu\nabla\widehat F(z))/\mu=\nabla F(x)-w/\mu$.
A bound on a full primal--dual product norm
bounds its primal component and hence every restricted primal chord.
Weak duality turns a feasible final duality gap at most $\eps_r$ into
primal accuracy. These hypotheses transfer the lower count; they do not
construct a dual-feasible version of the shorter primal route.

The box construction also yields the same lower counts on a matrix ball
or Jordan interval with at least $r$ channels, by placing the positive
weights in a fixed spectral frame and setting the other objective weights
to zero. Intermediate points need not stay in that frame. The contraction
in \cref{thm:exact-distance} bounds their distance to each labeled center,
and the spectral trace inequality turns final accuracy into a lower bound
$1-\eps_r$ on the largest singular value or Jordan eigenvalue. These are
exactly the two facts about arbitrary iterates used in the proof above.

\subsection{Unequal scales can produce a larger discrete separation}
The scale frontier also has a finite-sequence interpretation, with an
explicit parameter-crossing assumption distinct from \eqref{eq:tube-contract}.
In this subsection the barrier is
$F_\alpha(x)=\sum_{i=1}^r\alpha_i b(x_i)$ on $(-1,1)^r$;
all local norms, distances, lengths, and labeled centers refer to this
weighted barrier.
Fix $B\geq4$ and an integer $r$. For large $U$, put
\[
 u_i=UB^{r-i},\quad H_0=H(u_1),\quad
 \alpha_i=\left(\frac{H_0}{H(u_i)}\right)^2\geq1,\quad
 w_i=\alpha_i e^{u_i},\qquad i=1,\ldots,r.
\]
At terminal label zero all weighted radial coordinates equal $H_0$;
thus the endpoint distance is $H_0\sqrt r$. The central length is
$\Theta(H_0r)$, uniformly for large $U$: for fixed $M>0$ the disjoint
cores $I_i=[-u_i+M,-u_{i+1}-M]$, $i<r$, and
$I_r=[-u_r+M,0]$ each have gain $\Theta(H_0)$ in their $i$th weighted
radial coordinate. Indeed, $H(t)=t+O(1)$ for positive $t$ makes this gain
at least a fixed fraction of $H_0$; the upper bound follows from
\cref{thm:sharp-prefix} and $\Gamma\leq\sqrt r$.

For fixed $0<R<1$ and $\delta\geq0$, consider forward-$R$ iterates
starting at $z_0=0$
and within distance $\delta$ of their labeled centers, with
$s_0\leq-u_1+M$ and $s_N\geq0$; the initial label may be $-\infty$.
They need $\Omega_{R,\delta}(H_0r)$ moves once $U$ is large enough.
To see this, sum the clipped $i$th coordinate gain across each core as
a scalar potential. Every full core has gain exceeding
$C=2\delta-\log(1-R)$. A permitted center-to-center jump can meet at
most two cores in positive length: meeting three would cross the middle
one fully, contradicting its distance bound. The at most two clipped
coordinate gains sum to at most $\sqrt2C$ by Cauchy--Schwarz and exact
coordinate distance. Net progress is $\Theta(H_0r)$, even with
backtracking. To compare the same endpoints, specialize to
$z_N=x(0)$ and $s_N=0$. The lower bound still applies, while sampling
a shortest route from $0$ to $x(0)$ takes $O_R(H_0\sqrt r)$ moves.
This construction uses large scale and objective
dynamic ranges; \eqref{eq:scale-frontier} explains why such ranges are
needed. No ordinary-input efficiency claim is made for this construction.
